../cictikz/src/cictikz/data/tex/cictikz_preamble.tex

% Signal flow graph of a general biquad: the first-order graph with a second
% integrator added, so it uses the same summing node and 1/s block from
% sfg_lib.  V_i reaches the first sum through k_0/omega_0 and the second
% through k_1 and k_2 s; V_o is fed back through -omega_0 to the first sum
% and through -omega_0/Q to the second, giving
%   V_o/V_i = (k_2 s^2 + k_1 s + k_0)/(s^2 + (omega_0/Q) s + omega_0^2)

\begin{circuitikz}[american, thick, transform shape]

  %%%%%%%%%%%%%%%%%%%%%%%%%%%%%%%%%%%%%%%%%%%%%%%%%%%%%%%%%%%%%%%%%%%%%%
%% Shared pieces for the l04_afe signal flow graphs, l4_first_order and
%% l4_biquad.  The biquad is meant to read as the first-order graph with a
%% second integrator added, so the summing node and the 1/s block have to be
%% the same size and shape in both.  They are the same size in the original
%% artwork too.
%%
%% Names are prefixed and letters only, and computed coordinates are braced:
%% TikZ scans a coordinate for its closing parenthesis, so a '(' inside an
%% expression ends it early.
%%%%%%%%%%%%%%%%%%%%%%%%%%%%%%%%%%%%%%%%%%%%%%%%%%%%%%%%%%%%%%%%%%%%%%

\newcommand{\sfgR}{0.25}    % summing node radius
\newcommand{\sfgBw}{1.0}    % half width of the 1/s block
\newcommand{\sfgBh}{0.65}   % half height of the 1/s block

% A summing node at (#1,#2).
\newcommand{\sfgSum}[2]{%
  \draw (#1,#2) circle (\sfgR);
  \draw ({#1-0.17},#2) -- ({#1+0.17},#2);
  \draw (#1,{#2-0.17}) -- (#1,{#2+0.17});
}

% A block centred at (#1,#2) carrying #3.
\newcommand{\sfgBox}[3]{%
  \draw ({#1-\sfgBw},{#2-\sfgBh}) rectangle ({#1+\sfgBw},{#2+\sfgBh});
  \node at (#1,#2) {#3};
}

\newcommand{\sfgDot}[2]{\fill (#1,#2) circle (0.07);}


  \def\ym{2.7}    % the forward path
  \def\ykone{1.45}
  \def\yktwo{0.6}
  \def\yq{3.7}    % the -omega_0/Q path
  \def\yw{4.9}    % the -omega_0 path
  \def\xin{1.3}   % where the k_1 and k_2 s paths leave the input
  \def\xsa{3.5}   % first summing node
  \def\xsb{7.75}  % second summing node
  \def\xt{11.2}   % where the two feedbacks tap the output

  % forward path
  \node[anchor=east] at (0.25,\ym) {$V_{i}$};
  \draw[->] (0.35,\ym) -- ({\xsa-\sfgR},\ym);
  \sfgSum{\xsa}{\ym}
  \draw[->] ({\xsa+\sfgR},\ym) -- (4.5,\ym);
  \sfgBox{5.5}{\ym}{$\frac{1}{s}$}
  \draw[->] (6.5,\ym) -- ({\xsb-\sfgR},\ym);
  \sfgSum{\xsb}{\ym}
  \draw[->] ({\xsb+\sfgR},\ym) -- (8.75,\ym);
  \sfgBox{9.75}{\ym}{$\frac{1}{s}$}
  \draw[->] (10.75,\ym) -- (11.6,\ym);
  \node[anchor=west] at (11.7,\ym) {$V_{o}$};

  % k_1 and k_2 s, both taken off the input and summed into the second node
  \draw (\xin,\ym) -- (\xin,\yktwo);
  \sfgDot{\xin}{\ym}
  \sfgDot{\xin}{\ykone}
  \draw (\xin,\ykone) -- (6.5,\ykone);
  \draw[->] (6.5,\ykone) -- (7.57,2.52);
  \draw (\xin,\yktwo) -- (\xsb,\yktwo);
  \draw[->] (\xsb,\yktwo) -- (\xsb,{\ym-\sfgR});

  % the two feedbacks, both summed from above
  \draw (\xt,\ym) -- (\xt,\yw) -- (\xsa,\yw);
  \draw[->] (\xsa,\yw) -- (\xsa,{\ym+\sfgR});
  \draw (\xt,\yq) -- (\xsb,\yq);
  \draw[->] (\xsb,\yq) -- (\xsb,{\ym+\sfgR});
  \sfgDot{\xt}{\ym}
  \sfgDot{\xt}{\yq}

  \node[anchor=south] at (2.3,{\ym+0.1}) {$k_{0}/\omega_{0}$};
  \node[anchor=south] at (6.95,{\ym+0.1}) {$\omega_{0}$};
  \node[anchor=north] at (3.8,{\ykone-0.1}) {$k_{1}$};
  \node[anchor=north] at (4.25,{\yktwo-0.1}) {$k_{2}s$};
  \node[anchor=south] at (7.1,{\yw+0.1}) {$-\omega_{0}$};
  \node[anchor=south] at (9.85,{\yq+0.1}) {$-\omega_{0}/Q$};

\end{circuitikz}
\end{document}
